% Insert for the low-prime rank lemma.
% This file is intentionally separate from every frozen paper directory.

\subsection{The low-prime rank lemma}
\label{sec:insert-low-prime-rank}

This subsection gives the coefficientwise local argument for the low-prime
branch. It applies to the pencil $E+TC$, including finite parts at poles and
the coefficient of $T$.
Throughout this subsection
\[
 I=\{0,\ldots,m-1\},\qquad J=\{m,\ldots,2m-1\},\qquad U=2m-1,
\]
\[
 w(X)=\frac{\mathcal N(X)}d,\qquad
 \mathcal N(X)=\prod_{u=1}^{g}(X-u)^3,\qquad d=(g!)^3.
\]
Thus $w(r)=\binom{r-1}{g}^{3}$ at every non-negative integer $r$.

\begin{lemma}[Low-prime rank]
\label{lem:residue-copy-rank-insert}
Let $p\geq5$ satisfy
\[
 p^2>U+g,
 \qquad p\nmid a(a-1),
 \qquad q=\left\lfloor\frac gp\right\rfloor\geq1,
\]
where $a$ is an integer with $|a|>1$.  Write
\[
 g=qp+h,\qquad 0\leq h<p,\qquad G=p-h.
\]
Call a node $n$ good if $n\bmod p<G$ and bad if
$n\bmod p\geq G$.  Put
\[
 \kappa_n=\begin{cases}0,&n\text{ good},\\1,&n\text{ bad},\end{cases}
 \qquad
 \omega_i=2-\kappa_i,\qquad \omega'_j=2-\kappa_j.
\]
For $n\geq0$ define
\[
 R_n(X)=\frac{\mathcal N(X)-\mathcal N(-n)}{X+n},
 \qquad
 H_n(a)=\sum_{s=1}^{n}\frac{a^{n-s}}{s^3},
\]
and
\begin{equation}
 b_n(T)=\sum_{r\geq1}a^{-r}\left(\frac{R_n}{d}\right)''(r)
       +2w(-n)\bigl(a^nT-H_n(a)\bigr).
 \label{eq:insert-bn}
\end{equation}
Let the normalized pencil entries be
\begin{equation}
 B_{ij}(T)=\frac{b_i(T)-b_j(T)}{j-i},
 \qquad i\in I,\ j\in J.
 \label{eq:insert-Bij}
\end{equation}
By the second divided-difference identity and the finite harmonic terms,
these normalized entries are exactly the entries of the determinant pencil:
\[
 B_{ij}(T)=E_{ij}(a)+TC_{ij}(a).
\]
Then the exact row and column scaling is
\begin{equation}
 \widehat B_{ij}(T)=p^{\omega_i+\omega'_j}B_{ij}(T)
 =p^{(2-\kappa_i)+(2-\kappa_j)}B_{ij}(T).
 \label{eq:insert-scaled-B}
\end{equation}
The matrix $\widehat B$ has entries in $\mathbb Z_{(p)}[T]$, and, if
$N_I$ and $N_J$ denote the numbers of good nodes in $I$ and $J$, respectively,
\begin{equation}
 \operatorname{rank}_{\mathbb F_p(T)}\overline{\widehat B}
 \leq \min(N_I,N_J)+\min(G,h).
 \label{eq:insert-rank-bound}
\end{equation}
Here the bar denotes reduction in the Gauss valuation ring of
$\mathbb Q_p(T)$.  Consequently, with
\[
 W_{\rm nat}=\sum_{i\in I}\omega_i+\sum_{j\in J}\omega'_j
             =2m+N_I+N_J,
\]
and put
\[
 \varrho_*:=\min(N_I,N_J)+\min(G,h).
\]
every coefficient of the determinant satisfies
\begin{equation}
 v_p\!\left([T^k]\det B(T)\right)
 \geq -W_{\rm nat}+m-\varrho_*
 \qquad(0\leq k\leq m).
 \label{eq:insert-rank-valuation}
\end{equation}

There is also a separate entrywise alternative.  Its row and column base
weights are $1$ and $2$, respectively.  For each good residue $s<G$, add one
extra power of $p$ to the side having fewer copies of that residue.  If
\[
 c_I(s)=\#\{i\in I:i\equiv s\pmod p\},\qquad
 c_J(s)=\#\{j\in J:j\equiv s\pmod p\},
\]
the resulting total exponent is
\begin{equation}
 W_{\rm alt}=3m+\sum_{0\leq s<G}\min\{c_I(s),c_J(s)\}.
 \label{eq:insert-W-alt}
\end{equation}
Either alternative alone supplies an integerizer; their savings are not
combined.
\end{lemma}

\begin{proof}
The proof has three steps: summation, local order, and rank.

\paragraph{\emph{Newton summation modulo $p$.}}
We first need an elementary summation fact.  Suppose
\[
 V(X)=\sum_{\ell=0}^{D}v_\ell\binom X\ell,
 \qquad v_\ell\in\mathbb Z_{(p)}.
\]
The ordinary series $\sum_{r\geq1}a^{-r}V(r)$ converges because $|a|>1$.
We claim that it lies in $\mathbb Z_{(p)}$.  Write $z=a^{-1}$.  Then
\begin{equation}
 \sum_{r\geq1}z^r\binom r0=\frac z{1-z},
 \qquad
 \sum_{r\geq1}z^r\binom r\ell
 =\frac{z^\ell}{(1-z)^{\ell+1}}\quad(\ell\geq1).
 \label{eq:insert-newton-sums}
\end{equation}
The only denominators are powers of $1-a^{-1}$.  This is a $p$-unit because
$p\nmid a(a-1)$.  Formula \eqref{eq:insert-newton-sums} also shows that
reduction modulo $p$ may be done term by term in the finite Newton
expansion.  We use no $p$-adic convergence of the infinite series.

We also need a periodic form.  Let $V$ be a polynomial with
$p$-integral Newton coefficients.  Suppose that
$V(r+p^2)\equiv V(r)\pmod p$ for every integer $r$.  Then
\begin{equation}
 \sum_{r\geq1}a^{-r}V(r)
 \equiv
 \frac{\displaystyle\sum_{r=1}^{p^2}a^{-r}V(r)}
      {1-a^{-p^2}}\pmod p.
 \label{eq:insert-periodic-sum}
\end{equation}
To see this, write $S=\sum_{r\geq1}a^{-r}V(r)$ and
$V(r+p^2)=V(r)+pV_1(r)$.  The polynomial $V_1$ is integer valued at
$p$-adic integers.  So it has $p$-integral Newton coefficients.  By the
first fact its sum is $p$-integral.  Splitting off the first $p^2$ terms gives
$S=\sum_{r=1}^{p^2}a^{-r}V(r)+a^{-p^2}(S+pS_1)$ with $S_1$ $p$-integral.
The denominator $1-a^{-p^2}$ is a unit since
$a^{-p^2}\equiv a^{-1}\not\equiv1\pmod p$.

\paragraph{\emph{Local order.}}
For a residue class $t\pmod p$, let
\[
 \mathfrak m_t=(p,X-t)\subset\mathbb Z_{(p)}[X]
\]
and let $\operatorname{ord}_t(f)$ be the largest $e$ for which
$f\in\mathfrak m_t^e$.  The local ring at $\mathfrak m_t$ is regular of
dimension two.  Its associated graded ring is
\begin{equation}
 \operatorname{gr}_{\mathfrak m_t}\mathbb Z_{(p)}[X]
 \simeq\mathbb F_p[P,Y],
 \qquad P=\operatorname{in}(p),\quad Y=\operatorname{in}(X-t).
 \label{eq:insert-associated-graded}
\end{equation}
This graded ring is an integral domain.  Three rules follow at once:
\begin{enumerate}
\item exact division by $X-t$ lowers local order by one.  Division by
  $X-u$ with $u\not\equiv t\pmod p$ is division by a unit.  It does not
  lower local order;
\item differentiation lowers local order by at most one;
\item if $x\equiv t\pmod p$, then evaluation gives
  $v_p(f(x))\geq\operatorname{ord}_t(f)$.
\end{enumerate}
The second rule follows by differentiating monomials
$p^r(X-t)^s$.  The third follows by substituting
$X-t=p\,u$ with $u\in\mathbb Z_p$.

Since $g<p^2$, Legendre's formula gives
\begin{equation}
 v_p(d)=3v_p(g!)=3q.
 \label{eq:insert-d-valuation}
\end{equation}
Every residue class occurs $q$ times among $1,\ldots,g$.  The classes
$1,\ldots,h$ occur once more.  Therefore
\begin{equation}
 \operatorname{ord}_t\mathcal N\geq3q,
 \qquad
 \operatorname{ord}_t\mathcal N\geq3q+3
 \quad\text{if }t\equiv1,\ldots,h\pmod p.
 \label{eq:insert-N-orders}
\end{equation}
If $n$ is bad, then $-n\equiv1,\ldots,h\pmod p$.  So
$\mathcal N(-n)$ has valuation at least $3q+3$.  If $n$ is good, its
valuation is at least $3q$.

\paragraph{\emph{Scaled node sums.}}
Take a residue $t\not\equiv-n\pmod p$.  There the denominator $X+n$ is a
unit, and \eqref{eq:insert-N-orders} gives
\[
 \operatorname{ord}_tR_n\geq3q.
\]
At $t\equiv-n\pmod p$ division lowers order by one.  The numerator then
has stronger bounds.  For bad $n$ both terms in the numerator have order at
least $3q+3$.  Hence $\operatorname{ord}_tR_n\geq3q+2$.  For good $n$ the
bound is $\operatorname{ord}_tR_n\geq3q-1$.  We take two derivatives and
divide by $d$.  The local-order rules then give
\begin{equation}
 p^2\left(\frac{R_n}{d}\right)''(r)\in\mathbb Z_p
 \quad(n\text{ bad}),
 \qquad
 p^3\left(\frac{R_n}{d}\right)''(r)\in\mathbb Z_p
 \quad(n\text{ good})
 \label{eq:insert-Rn-values}
\end{equation}
for every integer $r$.  By the finite-difference characterization of
$\operatorname{Int}(\mathbb Z,\mathbb Z_p)$, the scaled polynomials have
$p$-integral Newton coefficients.  So the Newton summation fact applies to
their complete sums.

Now take the finite-part term.  By hypothesis $n<U+1<p^2$.  So
$v_p(s)\leq1$ for $1\leq s\leq n$, and
$v_p(H_n(a))\geq-3$.  Also
\begin{equation}
 v_p(w(-n))\geq3\quad(n\text{ bad}),
 \qquad v_p(w(-n))\geq0\quad(n\text{ good}).
 \label{eq:insert-harmonic-bounds}
\end{equation}
So the $T$ and harmonic terms in \eqref{eq:insert-bn} have the same
scaled integrality as in \eqref{eq:insert-Rn-values}.  We have proved
\begin{equation}
 p^2b_n(T)\in\mathbb Z_{(p)}[T]\quad(n\text{ bad}),
 \qquad
 p^3b_n(T)\in\mathbb Z_{(p)}[T]\quad(n\text{ good}).
 \label{eq:insert-bn-integrality}
\end{equation}

\paragraph{\emph{The one-pole reduction.}}
Suppose that $i$ is bad and $j$ is good.  Their residues are distinct.  So
$j-i$ is a $p$-unit.  The first congruence in
\eqref{eq:insert-bn-integrality} gives
$p^3b_i(T)\equiv0\pmod p$.  The coefficient of $T$ in $p^3b_j(T)$ is
also zero modulo $p$, because $p^3w(-j)a^j\equiv0\pmod p$.  Define
\begin{equation}
 \mathfrak d_j(a)=p^3b_j(0)\bmod p.
 \label{eq:insert-Dj}
\end{equation}
It depends on neither $i$ nor $T$.  Since $p^3b_i(T)\equiv0$,
\eqref{eq:insert-Bij} gives $p^3B_{ij}(T)\equiv-\mathfrak d_j(a)/(j-i)$.  That is,
\begin{equation}
 p^3B_{ij}(T)\equiv \frac{\mathfrak d_j(a)}{i-j}\pmod p
 \qquad(i\text{ bad},\ j\text{ good}).
 \label{eq:insert-one-pole}
\end{equation}
The transpose formula holds for a good row and a bad column.

\paragraph{\emph{The bad--bad reduction.}}
Let $i$ and $j$ both be bad.  Put
\begin{equation}
 \delta=j-i,\qquad
 L_{ij}(X)=
 \frac{\mathcal N(-i)(X+j)-\mathcal N(-j)(X+i)}{j-i}.
 \label{eq:insert-Lij}
\end{equation}
Then the exact polynomial identity
\begin{equation}
 \mathcal N(X)-L_{ij}(X)=(X+i)(X+j)Q_{ij}(X)
 \label{eq:insert-Qij}
\end{equation}
defines $Q_{ij}\in\mathbb Z_{(p)}[X]$.  Since $0<|\delta|<2m<p^2$, we have
$v_p(\delta)\leq1$.  The two endpoint values of $\mathcal N$ have
valuation at least $3q+3$.  So every coefficient of $L_{ij}$ satisfies
\begin{equation}
 v_p\bigl([X^\ell]L_{ij}\bigr)\geq3q+2.
 \label{eq:insert-Lij-valuation}
\end{equation}

At an evaluation residue the vanishing factors fall into three cases.
Write $t\equiv r\pmod p$.  First suppose that the opposite node residue
$s\equiv-r\pmod p$ is bad.  Then
$\operatorname{ord}_t\mathcal N\geq3q+3$.
\begin{enumerate}
\item If both $X+i$ and $X+j$ vanish at $t$, then $i\equiv j\pmod p$ and
  $v_p(\delta)=1$.  Use the coordinate $X=t+Y$.  The constant coefficient
 of $L_{ij}$ has valuation at least $3q+3$.  Here the extra factor $t+i$
  or $t+j$ supplies one $p$ before division by $\delta$.  The coefficient
  of $Y$ has valuation at least $3q+2$.  Hence
  $\operatorname{ord}_tL_{ij}\geq3q+3$.
  Dividing by two vanishing factors gives
  $\operatorname{ord}_tQ_{ij}\geq3q+1$.
\item If exactly one factor vanishes, the other is a unit.  From
  \eqref{eq:insert-Lij-valuation} and the order $3q+3$ of $\mathcal N$,
  division by the one vanishing factor gives
  $\operatorname{ord}_tQ_{ij}\geq3q+1$.
\item If neither factor vanishes, both are units.  Equations
  \eqref{eq:insert-N-orders} and \eqref{eq:insert-Lij-valuation} give
  $\operatorname{ord}_tQ_{ij}\geq3q+2$.
\end{enumerate}
Divide by $d$ and differentiate twice.  Then multiplication by $p^2$ makes
each of the three cases divisible by $p$.  Thus
\begin{equation}
 p^2\left(\frac{Q_{ij}}d\right)''(r)\equiv0\pmod p
 \quad\text{when }-r\text{ is a bad residue}.               
 \label{eq:insert-Q-bad-residue}
\end{equation}

Now let $s\equiv-r\pmod p$ be good.  Both $X+i$ and $X+j$ are units in the
local ring at $t=-s$.  Set
$u(X)=((X+i)(X+j))^{-1}$.  The derivatives $u,u',u''$ are integral there.
All coefficients of $L_{ij}/d$ are divisible by $p^2$.  So the $L_{ij}$ term
vanishes after the present reduction.  Consider the terms of $(wu)''$
containing $w$ and $w'$.  By the local-order bounds for $w,w',w''$, they also
vanish after multiplication by $p^2$.  Consequently
\begin{equation}
 p^2\left(\frac{Q_{ij}}d\right)''(r)
 \equiv\frac{p^2w''(r)}{(r+i)(r+j)}\pmod p
 \quad\text{when }-r\text{ is good}.                         
 \label{eq:insert-Q-good-residue}
\end{equation}

The identity \eqref{eq:insert-Qij} is exactly the divided-difference identity
\[
 \frac{R_i(X)-R_j(X)}{j-i}=Q_{ij}(X).
\]
Therefore \eqref{eq:insert-Bij} takes the following form:
\begin{multline}
 B_{ij}(T)=\sum_{r\geq1}a^{-r}\left(\frac{Q_{ij}}d\right)''(r)\\
 +\frac{2w(-i)(a^iT-H_i(a))-2w(-j)(a^jT-H_j(a))}{j-i}.
 \label{eq:insert-B-Q}
\end{multline}
For bad $i,j$ consider the second term in \eqref{eq:insert-B-Q}.  We have
$v_p(w(-i)),v_p(w(-j))\geq3$, $v_p(H_i),v_p(H_j)\geq-3$ and
$v_p(j-i)\leq1$.  So its $T$ part has valuation at least $2$.  Its
harmonic part has valuation at least $-1$.  After multiplication by $p^2$
both vanish modulo $p$.

It remains to sum the good-residue contribution.  Define
\[
 \varpi(r)=p^2w''(r)\pmod p.
\]
The polynomial $p^2w''$ has $p$-integral Newton coefficients.  Its values
have period $p^2$ modulo $p$.  To see this, apply Taylor expansion and the
local-order rule.  For every $\ell\geq1$ they give
\begin{equation}
 v_p\!\left(\frac{p^{2\ell+2}}{\ell!}
              w^{(\ell+2)}(r)\right)
 \geq \ell-v_p(\ell!)\geq1,
 \label{eq:insert-z-period}
\end{equation}
where the last inequality uses $p\geq5$.  Now take $V=p^2(Q_{ij}/d)''$.
By the case analysis above its values are $p$-integral.  So it has
$p$-integral Newton coefficients.  By
\eqref{eq:insert-Q-bad-residue} and \eqref{eq:insert-Q-good-residue},
$V(r)\equiv \varpi(r)/((r+i)(r+j))$ when $-r$ is good and $V(r)\equiv0$
otherwise.  The factor $1/((r+i)(r+j))$ modulo $p$ depends only on
$r\bmod p$.  So $V(r+p^2)\equiv V(r)\pmod p$.  Hence
\eqref{eq:insert-periodic-sum} applies to $V$.  For a good residue $s<G$, set
\begin{equation}
 \mathfrak c_s(a)=\frac{\displaystyle
       \sum_{\substack{1\leq r\leq p^2\\-r\equiv s\ (p)}}a^{-r}\varpi(r)}
       {1-a^{-p^2}}\pmod p.
 \label{eq:insert-Cs}
\end{equation}
Grouping the terms in \eqref{eq:insert-B-Q} by $s\equiv-r\pmod p$ gives
\begin{equation}
 p^2B_{ij}(T)
 \equiv \sum_{0\leq s<G}
       \frac{\mathfrak c_s(a)}{(i-s)(j-s)}\pmod p
 \qquad(i,j\text{ bad}).
 \label{eq:insert-bad-bad}
\end{equation}
All displayed denominators are $p$-units, since $i,j$ are bad and $s$ is
good.  The expression is independent of $T$.

\paragraph{\emph{Rank count and valuation transfer.}}
Take a bad row $i$ and a good column $j$.  Then the reduction
\eqref{eq:insert-one-pole} is a Cauchy feature with the good residue $s\equiv j\pmod p$.
For a bad row and a bad column, \eqref{eq:insert-bad-bad} is a sum of the
same features $1/(i-s)$, $s<G$.  Now take the block of bad rows, with all
columns.  It factors through the $G$ good residue features.  So its rank is
at most $G$.  Bad rows with the same residue modulo $p$ are also identical in
every column.  For good columns this follows from
\eqref{eq:insert-one-pole}.  For bad columns it follows from
\eqref{eq:insert-bad-bad}.  So the rank is also at most $h$, the number of
bad residue classes.  The bad-row block has rank at most $\min(G,h)$.
Adding the $N_I$ good rows gives
\begin{equation}
 \operatorname{rank}\overline{\widehat B}
 \leq N_I+\min(G,h).
 \label{eq:insert-rank-I}
\end{equation}
The same argument for the transpose gives the bound
$N_J+\min(G,h)$.  Taking the smaller of the two proves
\eqref{eq:insert-rank-bound}.

The factors in \eqref{eq:insert-scaled-B} are the row and column
weights.  Bad nodes receive one power and good nodes two powers.  Thus
\begin{equation}
 \det\widehat B=p^{W_{\rm nat}}\det B.
 \label{eq:insert-det-scaling}
\end{equation}
The Gauss valuation on $\mathbb Q_p(T)$ has residue field $\mathbb F_p(T)$.
Let a matrix have residue rank at most $r$.  Gaussian elimination over the
valuation ring shows that its determinant has valuation at least $m-r$.
Equivalently, every $T$-coefficient of the determinant has valuation at
least $m-r$.  We apply this to $\widehat B$.  With
\eqref{eq:insert-rank-bound} it proves \eqref{eq:insert-rank-valuation}.

Finally, the one-pole bounds give $p^3B_{ij}\in\mathbb Z_{(p)}[T]$ for
distinct residues.  A same-good-residue pair can lose one further power in
the divided difference.  It is covered by $p^4B_{ij}$.  The bad--bad and
mixed cancellations above show that $p^3B_{ij}$ is integral in all remaining
cases.  So one row power and two column powers cover every edge except the
same-good-residue edges.  For a good residue $s$ only those edges need one
extra power.  Put that extra power on the side with fewer copies.  This gives
\eqref{eq:insert-W-alt} and completes the proof.
\qedhere
\end{proof}
